\documentclass[anon,12pt]{clear2025} % Anonymized submission
% \documentclass[final,12pt]{clear2025} % Final submission
% \documentclass[final,12pt,noproceedings]{clear2025} % Final submission for those who opt out of the proceedings; they should use clear2025-extended-abstract.tex for their extended abstract

% --- PACKAGES ---
\newtheorem{observation}[theorem]{Observation}
\newtheorem{attempt}{Attempt}

\usepackage[capitalise]{cleveref}
\usepackage{amsfonts}
\usepackage{bbm}

\usepackage{tikz}
\usetikzlibrary{bayesnet}

\usetikzlibrary{shapes,decorations,arrows,calc,arrows.meta,fit,positioning}
\tikzset{
    -Latex,auto,node distance =0.5 cm and 0.5 cm,semithick,
    state/.style ={ellipse, draw, minimum width = 0.7 cm},
    point/.style = {circle, draw, inner sep=0.04cm,fill,node contents={}},
    bidirected/.style={Latex-Latex,dashed},
    el/.style = {inner sep=2pt, align=left, sloped}
}


\usepackage{url}
\usepackage{booktabs}
\usepackage{enumitem}

\tikzset{%
  zeroarrow/.style = {-stealth,dashed},
  onearrow/.style = {-stealth,solid},
  c/.style = {circle,draw,solid,minimum width=2em,
        minimum height=2em},
  r/.style = {rectangle,draw,solid,minimum width=2em,
        minimum height=2em}
}

\makeatletter
\DeclareRobustCommand{\varname}[1]{\begingroup\newmcodes@\mathit{#1}\endgroup}
\makeatother

\newcommand{\midlinewidth}{1.0pt}
\newcommand{\incmidlinewidth}{1.7pt}
\newcommand{\middist}{20pt}
\newcommand{\halfdist}{13pt}
\newcommand{\pn}{\ensuremath{\mathit{PN}}}
\newcommand{\ps}{\ensuremath{\mathit{PS}}}
\newcommand{\pns}{\ensuremath{\mathit{PNS}}}

\definecolor{petroil2} {RGB} {36, 165, 175}
\definecolor{gold2} {RGB} {255, 130, 0}

% # CODE SNIPPET STYLE (settings borrowed from Edward ICLR paper)
\usepackage{listings}
\usepackage[dvipsnames]{xcolor}

\definecolor{strings}{rgb}{.624,.251,.259}
\definecolor{keywords}{rgb}{.224,.451,.686}
\definecolor{comment}{rgb}{.322,.451,.322}

\lstdefinelanguage{python}{
  morekeywords={from, import, as, for, in, while, def, return, 
  =, +, -, /, *, lambda},
  % keywords=[2]{build_toy_dataset, neural_network, __init__},
  keywords=[3]{sample, param, module, Marginal, Posterior, Trace, Poutine, Distribution},
  morecomment=[l]{\#},
  morecomment=[s]{"""}{"""},
  morestring=[b]',
  morestring=[b]",
  alsoletter={<>=-+/*},
  sensitive=true
}

\usepackage[most]{tcolorbox}
\newtcolorbox{docode}[1][]{
    colback=gray!20,        % background color
    colframe=black,         % border color
    width=0.9\textwidth,    % box width
    arc=4mm,                % rounded corners
    boxrule=0.5pt,          % border thickness
    left=2em,               % left padding
    enhanced,
    fontupper=\ttfamily,    % font inside the box
    #1                      % allow optional overrides
}



\lstset{
  language=python,
  keywordstyle=\color{BrickRed}\bfseries\ttfamily,
  keywordstyle=[2]\color{Violet}\ttfamily,
  keywordstyle=[3]\color{keywords}\ttfamily,
  commentstyle=\color{comment}\ttfamily,
  stringstyle=\color{strings}\ttfamily,
  basicstyle=\fontsize{8pt}{8.25pt}\selectfont\ttfamily,
  basewidth=0.5em,
  % numbers=left,
  % numberstyle=\tiny,
  % stepnumber=1,
  columns=fixed,
  xleftmargin=2ex,%0.25in,
  % firstnumber=1,
  showstringspaces=false,
  mathescape=true,
  keepspaces=True,
  tabsize=2
}
\renewcommand{\texttt}[1]{\lstinline[basicstyle=\fontsize{8pt}{8.25pt}\selectfont\ttfamily]{#1}}


\usepackage{xcolor}   % for colors
\newcommand{\CITE}{\colorbox{green!30}{\textbf{CITE}}}
\newcommand{\TODO}{\colorbox{orange!30}{\textbf{TODO}}}
\newcommand{\ourapproach}{our approach}
\usepackage{hyperref}

% Configure hyperlinks
\hypersetup{
    colorlinks=true,      % use colored text for links
    linkcolor=blue,       % color of internal links (sections, TOC)
    urlcolor=blue,        % color of URLs
    citecolor=blue,       % color of citations
    filecolor=blue        % color of file links
}

\title[Explainable]{A Computationally Feasible Framework for  Causal Probabilistic Explanation}
\usepackage{times}
% Use \Name{Author Name} to specify the name.
% If the surname contains spaces, enclose the surname
% in braces, e.g. \Name{John {Smith Jones}} similarly
% if the name has a "von" part, e.g \Name{Jane {de Winter}}.
% If the first letter in the forenames is a diacritic
% enclose the diacritic in braces, e.g. \Name{{\'E}louise Smith}

% Two authors with the same address
% \clearauthor{\Name{Author Name1} \Email{abc@sample.com}\and
%  \Name{Author Name2} \Email{xyz@sample.com}\\
%  \addr Address}

% Three or more authors with the same address:
% \clearauthor{\Name{Author Name1} \Email{an1@sample.com}\\
%  \Name{Author Name2} \Email{an2@sample.com}\\
%  \Name{Author Name3} \Email{an3@sample.com}\\
%  \addr Address}

% Authors with different addresses:
\clearauthor{%
 \Name{Rafal Urbaniak} \Email{rafal@basis.ai} \\
 \Name{Eli Bingham} \Email{eli@basis.ai}\\
 \addr{Basis Research Institute}\\
 \Name{Poorva Garg} \Email{poorvagarg@cs.ucla.edu}\\
 \addr{University of California, Los Angeles} \\
 \Name{Jack Feser} \Email{jack@basis.ai} \\
 \Name{Daniel Waxman} \Email{dan@basis.ai} \\
 \addr{Basis Research Institute}
}

\begin{document}

  
\begin{definition}[Probabilistic Structural Causal Model]
A probabilistic structural causal model is a tuple
\[
M = \langle \mathbf{U}, \mathbf{V}, \mathbf{F}, P(\mathbf{u}) \rangle,
\]
where \(\mathbf{V} = \{V_1,\dots,V_n\}\) is a finite set of endogenous variables,
\(\mathbf{U} = \{U_1,\dots,U_n\}\) is a set of exogenous variables,
and \(\mathbf{F} = \{f_1,\dots,f_n\}\) is a set of structural assignments
\[
V_i := f_i(\mathbf{Pa}_i, U_i), \qquad i = 1,\dots,n,
\]
with \(\mathbf{Pa}_i \subseteq \mathbf{V}\setminus\{V_i\}\) denoting the parents of \(V_i\).
The distribution \(P(\mathbf{u})\) is a joint distribution over \(\mathbf{U}\),
and the assignments in \(\mathbf{F}\) are assumed to have a unique solution for \(\mathbf{V}\) given \(\mathbf{U}\).
\end{definition}




\begin{definition}[Causal set]
    Given a probabilistic causal model,  \(\langle \mathbf{U}, \mathbf{V}, \mathbf{F}, P(\mathbf{u})\rangle\), a causal set is a set of variable value pairs \(\{\mathbf{X} = \mathbf{x} ~|~ \mathbf{X} \subseteq  \mathbf{V}, \mathbf{x} \in \mathit{dom}(\mathbf{X})\}\). 
\end{definition}


\begin{definition}[Interventions]
Given a probabilistic structural causal model
\(M = \langle \mathbf{U}, \mathbf{V}, \mathbf{F}, P(\mathbf{u}) \rangle\),
an intervention on a set of variables \(\mathbf{X} = \{X_1, \dots, X_k\} \subseteq \mathbf{V}\)
with values \(\mathbf{x} = (x_1, \dots, x_k)\) replaces the structural assignments for \(X_i\) by $X_i := x_i$, where $i = 1, \dots, k$.  We denote this intervention by \([\mathbf{X} \leftarrow \mathbf{x}]\).
\end{definition}


\begin{definition}[Variable Selection Distribution]
Let $\mathbf{X}$ be a finite set of random variables and $2^\mathbf{X}$ be its power set. The random variable $\mathcal{Y}$ is distributed according to $\Gamma(2^\mathbf{X})$ such that $ \mathbf{X} \supseteq \mathbf{Y} \thicksim \Gamma(2^\mathbf{X})$ and $p_\Gamma(\mathbf{Y}) = \mathbb{P}(\mathcal{Y} = \mathbf{Y})$.
\end{definition}


\begin{example}[Cardinality-Constrained Uniform Selection]
Let $\mathbf{X}$ be a candidate set of random variables and let $k \thicksim \mathrm{Unif}(J, K)$ with $1 \leq J < K \leq |\mathbf{X}|$. A variable selection $\mathbf{Y} \thicksim \Gamma(2^\mathbf{X})$ is sampled from the cardinality-constrained uniform selection distribution $\Gamma(2^\mathbf{X})$ with
\[
p_\Gamma(\mathbf{Y}) \propto \mathbb{I}[|\mathbf{Y}| = k].
\]
\end{example}


\begin{definition}[Alternative Value Distribution]
Consider a set of random variables $\mathbf{X}$ and event $\mathbf{X}=\mathbf{x}\in dom(\mathbf{X})$. An alternative value $\mathbf{x}' \thicksim \Delta(\mathbf{x})$ is sampled from the alternative value distribution $\Delta(\mathbf{x})$, where $\mathbf{x}' \in dom (\mathbf{X})$ and $p_\Delta(\mathbf{x}') = \mathbb{P}(\mathbf{X}'=\mathbf{x'})$ with $dom(\mathbf{X}) = dom(\mathbf{X}')$.
\end{definition}

\begin{example}[Epsilon Excised Distribution]
Let factual event $\mathbf{X}=\mathbf{x}$ have probability $p(\mathbf{x})$.
The \emph{$\epsilon$-excised alternative value distribution} $\Delta_{\epsilon}(\mathbf{x})$
 where $\mathbf{x}' \sim \Delta(\mathbf{x})$ with probability:
\[
p_\Delta(\mathbf{x}') \propto 
\begin{cases}
    0, & \text{if } x' \in [\,x - \epsilon,\; x + \epsilon\,], \\[6pt]
    p(x'), & \text{otherwise}.
\end{cases}
\]
That is, probability mass in an $\epsilon$-neighborhood of the factual value $x$
is removed, and the remaining mass is renormalized.
\end{example}


\begin{definition}[search densities] Suppose we start with some active suspects and witnesses distributions, together with a distribution for alternative values, with corresponding probability functions $p^\Gamma_s(\cdot), p^\Gamma_w(\cdot)$ and $p_\Delta(\cdot)$. Then the \textbf{sufficiency intervention} is obtained by setting the values of $\mathbf{C}, \mathbf{T}$ to their factual values, $\mathbf{c}^\star, \mathbf{t}^\star$, while the \textbf{necessity intervention} is obtained by setting the values of $\mathbf{C}$ to $\mathbf{c'}$ and keeping the values of $\mathbf{T}$ fixed at $\mathbf{t}^\star$; these result in the corresponding intervened densities in two interventional regimes (where $\mathbf{u}$ is the noise), which can further be weighted by the corresponding activation probabilities, and marginalized over all witness sets and all causal candidate sets containing the variable of interest $X_k$. 


\begin{align}
{p}^s_k(y \vert \mathbf{u})    &= \sum_{\mathbf{C} \subseteq \mathbf{S}, \mathbf{T}\subseteq \mathbf{W}} \mathbb{I}[X_k \in \mathbf{C}]
p (y \vert \mathbf{u}, do(\mathbf{C} = \mathbf{c^\star}, \mathbf{T} = \mathbf{t^\star})) p^\Gamma_s(\mathbf{C}) p^\Gamma_w(\mathbf{T})
%p^{es}_{C,T}(y \vert \mathbf{u}, C, T) 
\\ 
{p}^n_k(y \vert \mathbf{u})   & = \sum_{\mathbf{C} \subseteq \mathbf{S}, \mathbf{T}\subseteq \mathbf{W}} \mathbb{I}[X_k \in \mathbf{C}]
\int 
p(y \vert \mathbf{u}, do(\mathbf{C} = \mathbf{c'}, \mathbf{T} = \mathbf{t^\star}))
p^\Gamma_s(\mathbf{C}) p^\Gamma_w(\mathbf{T}) p_\Delta(\mathbf{c}') d\, \mathbf{c}'
%p^{en}_{C, T}( y \vert \mathbf{u}, C, T) 
\end{align}
\end{definition}





\begin{definition}[Causal impact and its push-forward]
Let $y^s$ denote the outcome in the \emph{sufficiency world}, $y^n$ denote the outcome in the \emph{necessity world}, and $y^\star$ the outcome in the factual world. By construction, these are independent conditional on $\mathbf{u}$, and so we get the joint:
\begin{align}
 p(\langle y^s, y^n \rangle \mid \mathbf{u})  & = p^s_k(y^s \mid \mathbf{u}) \; {p}^n_k(y^n \mid \mathbf{u}).
\end{align}

\noindent which also generates a push-forward on any causal impact scoring function $ci(y^s, y^n, y^\star)$, whose expectation is a summary of causal impact.
\begin{align}\label{eq:score}
\mathbb{E}_{\mathbf{u}}\Big[\, \mathbb{E}_{y^s, y^n \mid \mathbf{u}} \big[ ci(y^s, y^n, y^\star) \big] \,\Big]
= \\  =\int 
   \int \int ci(y^s, y^n, y^\star) \; p^s_k(y^s \mid \mathbf{u}) \; p^n_k(y^n \mid \mathbf{u}) p(\mathbf{u}) 
   \, dy^s \, dy^n \, d\mathbf{u}. \nonumber
\end{align}
\end{definition}


\begin{example}[Absolute difference impact score] \label{ex:absolute_score}
One straightforward function of that sort would equally reward distance from the factual value $y^\star$ in the necessity world, and proximity to the factual value in the sufficiency world. Using the two in the scoring is an echo of Pearl's intuition to use the probability of both sufficiency and necessity.
\begin{align*}
ci(y^s, y^n, y^\star) & = \vert y^n - y^\star \vert   - \vert y^s - y ^\star \vert. 
\end{align*}
\end{example}






\newpage


The notions we proposed are not meant to explicate Halpern's notion of actual causality; rather, the goal was to combine the intuitions underlying Halpern's notion with those motivating Pearl's notion of probability of necessity and sufficiency to provide a more general tool for causal attribution. Nevertheless, to argue that this generalization at least preserves some of the spirit of Halpern's definition, in this section, we discuss the connection between actual causality and causal impact. Let us work with structural causal models whose all variables are binary or categorical.  In this context, we can use Halpern's modified definition of actual causality.

\begin{definition} \label{def:ac}
    \(\mathbf{C} = \mathbf{c}\) is an actual cause of \(\varphi\) in the causal setting \((M, \mathbf{u})\) if the following three conditions hold:
    \begin{itemize}
        \item \textbf{Factivity}: \((M, \mathbf{u}) \vDash (\mathbf{C} = \mathbf{c}\,)\) and \((M, \mathbf{u}) \vDash \varphi\),
        \item \textbf{Context-sensitive Necessity}: There is a set of variables \(\mathbf{T} \subseteq \mathbf{V}\) and an alternative setting \(\mathbf{c'}\) of variables in \(\mathbf{C}\) such that if \((M, \mathbf{u}) \vDash \mathbf{T} = \mathbf{t^\star}\), then
            \[(M, \mathbf{u}) \vDash [\mathbf{C} \leftarrow \mathbf{c'}, \mathbf{T} \leftarrow \mathbf{t^\star}] \neg \varphi,\]

\vspace{-3mm}
    
        \item \textbf{Minimality}: \(\mathbf{C}\) is minimal; there is no strict subset \(\mathbf{C_{sub}} \subset \mathbf{C}\) such that \(\mathbf{C_{sub}} = \mathbf{c_{sub}}\) satisfies the above two conditions where $\mathbf{c_{sub}}$ is the restriction of \(\mathbf{c}\) to the variables in \(\mathbf{C_{sub}}\). 
    \end{itemize}
\end{definition}


\begin{theorem} \label{th:ac->positive} Suppose (i) Variable selection distributions are full support, i.e.: $\forall \,\mathbf{C} \subseteq \mathbf{V}\,\,\,\,\,\, p^\Gamma_s(\mathbf{C}) > 0$ and 
   $\forall \,\mathbf{T} \subseteq \mathbf{V}\,\,\,\,\,\,  p^\Gamma_w(\mathbf{T}) > 0$,     (ii) Alternative value distribution has full support on non-factive values, i.e. $p_\Delta(\mathbf{c}') > 0 \text{ whenever } c'\neq c^\star$, (iii) We let $y = \mathbb{I}(\varphi)$  and use a causal impact score such that  $ci(y^s, y^n, y^\star) = \mathbb{I}(y^n \neq y^\star)$. Given the above assumptions:

If \,\, \(\mathbf{C} = \mathbf{c}\) is an actual cause of \(\varphi\) in the causal setting \((M, \mathbf{u})\), \,\, then for any $X_k\in \mathbf{C}$, we have  $\mathbb{E}_{y^s, y^n \mid \mathbf{u}} \big[ ci(y^s, y^n, y^\star) \big] > 0$.
\end{theorem}


\begin{proof} By the necessity condition in the definition of actual causality, there is a set of variables \(\mathbf{T} \subseteq \mathbf{V}\), factual settings $\mathbf{c^\star, t^\star}$  and an alternative setting \(\mathbf{c'}\) such that \((M, \mathbf{u}) \vDash (\mathbf{C} = \mathbf{c}\,)\) and \((M, \mathbf{u}) \vDash \varphi\), and   $(M, \mathbf{u}) \vDash [\mathbf{C} \leftarrow \mathbf{c'}, \mathbf{T} \leftarrow \mathbf{t^\star}] \neg \varphi$.

Let's start with working out why $p^s_k(y\vert u)>0$ when $y=y^\star =1$. Consider:

\begin{align*}
\underbrace{\mathbb{I}[X_k \in \mathbf{C}]}_{>0 \text{by assumption}} \underbrace{p (y \vert \mathbf{u}, do(\mathbf{C} = \mathbf{c^\star}, \mathbf{T} = \mathbf{t^\star}))}_{>0 \text{ by factivity}} \underbrace{p^\Gamma_s(\mathbf{C}) p^\Gamma_w(\mathbf{T})}_{>0 \text{ by (i)}}
\end{align*}
\noindent so:

\begin{align*}
{p}^s_k(y \vert \mathbf{u})    &= \underbrace{\sum_{\mathbf{C} \subseteq \mathbf{S}, \mathbf{T}\subseteq \mathbf{W}} \mathbb{I}[X_k \in \mathbf{C}]
p (y \vert \mathbf{u}, do(\mathbf{C} = \mathbf{c^\star}, \mathbf{T} = \mathbf{t^\star})) p^\Gamma_s(\mathbf{C}) p^\Gamma_w(\mathbf{T})}_{>0 \text{ as at least one summand is $>0$ and none can be negative}}
\end{align*}



similarly, let us see why $p^n_k(y\vert \mathbf{u}) >0$ if $y = 0$. The summation step is the same, and some of the multiplicands are positive by the same arguments, we only need to notice that (1) we have go one more summation step deeper
\begin{align*} 
&\underbrace{p(y \vert \mathbf{u}, do(\mathbf{C} = \mathbf{c'}, \mathbf{T} = \mathbf{t^\star}))}_{> 0 \text{ by the necessity assumption}}
\underbrace{p^\Gamma_s(\mathbf{C}) p^\Gamma_w(\mathbf{T}) p_\Delta(\mathbf{c}')}_{>0 \text{ as before }} \\ &
\underbrace{\mathbb{I}[X_k \in \mathbf{C}]}_{>0 \text{ as before }}
\underbrace{\int 
p(y \vert \mathbf{u}, do(\mathbf{C} = \mathbf{c'}, \mathbf{T} = \mathbf{t^\star}))
p^\Gamma_s(\mathbf{C}) p^\Gamma_w(\mathbf{T}) p_\Delta(\mathbf{c}') d\, \mathbf{c}'}_{>0 \text { as at least one summand $>0$ and none can be negative}} \\
{p}^n_k(y \vert \mathbf{u})   & = \underbrace{\sum_{\mathbf{C} \subseteq \mathbf{S}, \mathbf{T}\subseteq \mathbf{W}} \mathbb{I}[X_k \in \mathbf{C}]
\int 
p(y \vert \mathbf{u}, do(\mathbf{C} = \mathbf{c'}, \mathbf{T} = \mathbf{t^\star}))
p^\Gamma_s(\mathbf{C}) p^\Gamma_w(\mathbf{T}) p_\Delta(\mathbf{c}') d\, \mathbf{c}'}_{>0 \text{ as at least one summand $>0$ and none can be negative}}
\end{align*}

Now let's look at the positivity of the expectation. For any setting of $y^s$ and $y^n$ we know \eqref{eq:inexpectation} is non-negative by construction. But in particular for $y^s = 1$ and $y^n=0$ we know it's positive:
\begin{align}\label{eq:inexpectation}
\underbrace{ci(y^s, y^n, y^\star)}_{\text{=1 by definition of ci}} \; \underbrace{p^s_k(y^s \mid \mathbf{u}) \; p^n_k(y^n \mid \mathbf{u})}_{>0 \text{ by previous arguments } } \underbrace{p(\mathbf{u})}_{>0 \text{ by factivity }}\end{align}
\noindent and therefore, by the properties of summation:
\begin{align}\label{eq:exp}
\mathbb{E}_{y^s, y^n \mid \mathbf{u}} \big[ci(y^s, y^n, y^\star) \big]
=  \int 
   \int \int ci(y^s, y^n, y^\star) \; p^s_k(y^s \mid \mathbf{u}) \; p^n_k(y^n \mid \mathbf{u}) p(\mathbf{u}) 
   \, dy^s \, dy^n  
   \end{align}
\noindent which completes the argument.
\end{proof}


Moreover, the partial implication in the opposite direction also holds:


\begin{theorem} Suppose the assumptions of Theorem \ref{th:ac->positive} hold. If $y^\star = 1$ and $\mathbb{E}_{y^s, y^n \mid \mathbf{u}} \big[ci(y^s, y^n, y^\star) \big] >0$, there is a set of suspects containing $X_k$ for which the necessity condition of the definition of actual causality holds for the respective model.
\end{theorem}


\begin{proof} Inspect \eqref{eq:exp}: the right-hand side needs to be positive, so for some combination of $y^s, y^n$ we need to have:
\begin{align*}
ci(y^s, y^n, y^\star) \; p^s_k(y^s \mid \mathbf{u}) \; p^n_k(y^n \mid \mathbf{u}) p(\mathbf{u})  > 0
\end{align*}
but we also know that $y^\star =1$, so the only combination of $y^s, y^n$ for which $ci(y^s, y^n, y^\star) \neq 0$ is $y^s =1, y^n = 0$. So we have:
\begin{align}
    p^n_k(0 \mid \mathbf{u}) p(\mathbf{u})&  >0\\
    \sum_{\mathbf{C} \subseteq \mathbf{S}, \mathbf{T}\subseteq \mathbf{W}} \mathbb{I}[X_k \in \mathbf{C}]
\int 
p(y \vert \mathbf{u}, do(\mathbf{C} = \mathbf{c'}, \mathbf{T} = \mathbf{t^\star}))
p^\Gamma_s(\mathbf{C}) p^\Gamma_w(\mathbf{T}) p_\Delta(\mathbf{c}') d\, \mathbf{c}' & > 0
\end{align}
\noindent so there have to be some $\mathbf{C, T, c'}$ such that:
\begin{align}
p(0 \vert \mathbf{u}, do(\mathbf{C} = \mathbf{c'}, \mathbf{T} = \mathbf{t^\star})) > 0
\end{align}
\noindent Since all the noise is exogenous and fixed, the consequences of interventions are deterministic, and therefore  $(M, \mathbf{u}) \vDash [\mathbf{C} \leftarrow \mathbf{c'}, \mathbf{T} \leftarrow \mathbf{t^\star}] \neg \varphi$, which completes the argument.
\end{proof}






\end{document}